\section{Prelude}

The prelude presents the fixed set-theoretic foundation used by every Felix
theory.  Its unproved propositions are exact, confined views of compiled
foundation principles; the later construction of the natural numbers is an
ordinary checked development from those principles.

% Only the confined prelude builder may authorize the foundation presentations
% below.  This source is not an ordinary axiom module.

\subsection{Universe principles}

\begin{proposition}[Universe containment]\label{universe_contains}
    For all $n$ we have $n\in\cumul{n}$.
\end{proposition}

\begin{proposition}[Universe transitivity]\label{universe_transitive}
    For all $n,a$ we have if $a\in\cumul{n}$, then
    for all $x$ we have if $x\in a$, then $x\in\cumul{n}$.
\end{proposition}

\begin{proposition}[Closure under family unions]%
\label{universe_union_closed}
    For all $n,a$ we have if $a\in\cumul{n}$, then
    $\unions{a}\in\cumul{n}$.
\end{proposition}

\begin{proposition}[Closure under powersets]\label{universe_power_closed}
    For all $n,a$ we have if $a\in\cumul{n}$, then
    $\pow{a}\in\cumul{n}$.
\end{proposition}

\begin{proposition}[Closure under replacement]%
\label{universe_replacement_closed}
    $\forall n:\Set.\ \forall A:\Set.\ \forall f:\Set\to\Set.
        A\in\cumul{n}\implies
        (\forall x:\Set.\ x\in A\implies
            \Apply{f}{x}\in\cumul{n})\implies
        \repl{A}{f}\in\cumul{n}$.
\end{proposition}

\begin{proposition}[Universe minimality]\label{universe_minimal}
    $\forall n:\Set.\ \forall U:\Set.
        n\in U\implies
        (\forall A:\Set.\ A\in U\implies
            (\forall x:\Set.\ x\in A\implies x\in U))\implies
        (\forall A:\Set.\ A\in U\implies \unions{A}\in U)\implies
        (\forall A:\Set.\ A\in U\implies \pow{A}\in U)\implies
        (\forall A:\Set.\ \forall f:\Set\to\Set.
            A\in U\implies
            (\forall x:\Set.\ x\in A\implies \Apply{f}{x}\in U)\implies
            \repl{A}{f}\in U)\implies
        (\forall x:\Set.\ x\in\cumul{n}\implies x\in U)$.
\end{proposition}

\begin{proposition}[Nested universes]\label{universe_nested}
    For all $m,n$ we have if $n\in\cumul{m}$, then
    for all $x$ we have if $x\in\cumul{n}$, then $x\in\cumul{m}$.
\end{proposition}
\begin{proof}
    Follows by \auto{universe_transitive,universe_union_closed,universe_power_closed,universe_replacement_closed,universe_minimal}.
\end{proof}

\subsection{Least fixed points}

The bounded least-fixed-point principles genuinely quantify over logical
functions and predicates, so their exact typed form is retained.

\begin{proposition}[Least-fixed-point bound]\label{lfp_bound}
    $\forall D:\Set.\ \forall F:\Set\to\Set.
        \forall x:\Set.\ x\in\lfp{D}{F}\implies x\in D$.
\end{proposition}

\begin{proposition}[Least pre-fixed point]\label{lfp_least}
    $\forall D:\Set.\ \forall F:\Set\to\Set.\ \forall C:\Set.
        (\forall x:\Set.\ x\in\Apply{F}{C}\implies x\in C)\implies
        (\forall x:\Set.\ x\in C\implies x\in D)\implies
        (\forall x:\Set.\ x\in\lfp{D}{F}\implies x\in C)$.
\end{proposition}

\begin{proposition}[Least fixed point]\label{lfp_fixed}
    $\forall D:\Set.\ \forall F:\Set\to\Set.
        ((\forall x:\Set.\ x\in\Apply{F}{D}\implies x\in D)\land
        (\forall X:\Set.\ \forall Y:\Set.
            ((\forall x:\Set.\ x\in X\implies x\in Y)\land
             (\forall x:\Set.\ x\in Y\implies x\in D))\implies
            (\forall x:\Set.
                x\in\Apply{F}{X}\implies x\in\Apply{F}{Y})))\implies
        \lfp{D}{F}=\Apply{F}{\lfp{D}{F}}$.
\end{proposition}

\begin{proposition}[Least-fixed-point induction]\label{lfp_induct}
    $\forall D:\Set.\ \forall F:\Set\to\Set.
     \forall P:\Set\to\Prop.\ \forall a:\Set.
        ((\forall x:\Set.\ x\in\Apply{F}{D}\implies x\in D)\land
        (\forall X:\Set.\ \forall Y:\Set.
            ((\forall x:\Set.\ x\in X\implies x\in Y)\land
             (\forall x:\Set.\ x\in Y\implies x\in D))\implies
            (\forall x:\Set.
                x\in\Apply{F}{X}\implies x\in\Apply{F}{Y})))\implies
        a\in\lfp{D}{F}\implies
        (\forall z:\Set.
            z\in\Apply{F}{\sep{\lfp{D}{F}}{P}}
                \implies \Apply{P}{z})\implies
        \Apply{P}{a}$.
\end{proposition}

\subsection{Elementary set constructions}

The following propositions give the public first-order interface to global
choice and the elementary set constructors.

\begin{proposition}[Global choice]\label{choice_member}
    For all $A,x$ we have if $x\in A$, then $\choice{A}\in A$.
\end{proposition}

\begin{proposition}[Set extensionality]\label{setext}
    For all $A,B$ we have
    if for all $x$ if $x\in A$, then $x\in B$, then
    if for all $x$ if $x\in B$, then $x\in A$, then
    $A=B$.
\end{proposition}

\begin{proposition}[Empty set]\label{emptyset}
    For all $x$ we have $x\in\emptyset$ iff $\bot$.
\end{proposition}

\begin{proposition}[Unordered pairs]\label{pairset_iff}
    For all $a,b,x$ we have
    $x\in\upair{a}{b}$ iff $x=a$ or $x=b$.
\end{proposition}

\begin{proposition}[Ordered-pair equality]\label{pair_eq_iff}
    For all $a,b,x,y$ we have
    $(a,b)=(x,y)$ iff $a=x$ and $b=y$.
\end{proposition}

\begin{proposition}[First projection]\label{fst_eq}
    For all $a,b$ we have $\fst{(a,b)}=a$.
\end{proposition}

\begin{proposition}[Second projection]\label{snd_eq}
    For all $a,b$ we have $\snd{(a,b)}=b$.
\end{proposition}

\begin{proposition}[Powersets]\label{pow_iff}
    For all $A,B$ we have $B\in\pow{A}$ iff
    for all $x$ if $x\in B$, then $x\in A$.
\end{proposition}

\begin{proposition}[Family unions]\label{unions_iff}
    For all $A,x$ we have $x\in\unions{A}$ iff
    there exists $B\in A$ such that $x\in B$.
\end{proposition}

\subsection{Natural numbers}

We now derive successor, inductive sets, and the natural numbers inside the
least universe containing the empty set.

\begin{definition}[Successor]\label{suc}
    $\suc{x} = \unions{\upair{\upair{x}{x}}{x}}$.
\end{definition}

\begin{definition}[Transitive sets]\label{transitiveset}
    $A$ is \in-transitive iff
    for every $x\in A$ we have
    for every $y\in x$ we have $y\in A$.
\end{definition}

\begin{definition}[Inductive sets]\label{inductive_set}
    $A$ is an inductive set iff
    $\emptyset\in A$ and
    for every $x\in A$ we have $\suc{x}\in A$.
\end{definition}

\begin{definition}[Ambient universe]\label{omega_universe}
    $\omegaUniverse = \cumul{\emptyset}$.
\end{definition}

\begin{proposition}[The empty set lies in the ambient universe]%
\label{omega_universe_empty}
    $\emptyset \in \omegaUniverse$.
\end{proposition}
\begin{proof}
    Follows by \auto{universe_contains,omega_universe}.
\end{proof}

\begin{proposition}[The empty set is transitive]\label{emptyset_transitive}
    $\emptyset$ is \in-transitive.
\end{proposition}
\begin{proof}
    Follows by \auto{transitiveset}.
\end{proof}

\begin{proposition}[Successors preserve transitivity]%
\label{successor_transitive}
    For all $x$ we have if $x$ is \in-transitive, then
    $\suc{x}$ is \in-transitive.
\end{proposition}
\begin{proof}
    Follows by \auto{transitiveset,suc}.
\end{proof}

\begin{proposition}[Successor powerset bound]\label{successor_power_bound}
    For all $x$ we have if $x$ is \in-transitive, then
    $\suc{x}\in\pow{\pow{x}}$.
\end{proposition}
\begin{proof}
    Follows by \auto{transitiveset,suc}.
\end{proof}

\begin{proposition}[Ambient successor closure]%
\label{omega_universe_transitive_successor}
    For all $x$ we have if $x\in\omegaUniverse$, then
    if $x$ is \in-transitive, then $\suc{x}\in\omegaUniverse$.
\end{proposition}
\begin{proof}
    Follows by \auto{universe_transitive,universe_power_closed,successor_power_bound,omega_universe}.
\end{proof}

\begin{definition}[Transitive part of the ambient universe]%
\label{omega_transitive_part}
    $\omegaTransitivePart = \{ x \in \omegaUniverse \mid \text{$x$ is \in-transitive} \}$.
\end{definition}

\begin{proposition}[The transitive part contains the empty set]%
\label{omega_transitive_part_empty}
    $\emptyset \in \omegaTransitivePart$.
\end{proposition}
\begin{proof}
    Follows by \auto{omega_transitive_part,omega_universe_empty,emptyset_transitive}.
\end{proof}

\begin{proposition}[The transitive part is successor-closed]%
\label{omega_transitive_part_successor}
    For all $x$ we have if $x\in\omegaTransitivePart$, then
    $\suc{x}\in\omegaTransitivePart$.
\end{proposition}
\begin{proof}
    Follows by \auto{omega_transitive_part,omega_universe_transitive_successor,successor_transitive}.
\end{proof}

\begin{proposition}[Infinity]\label{infinity}
    $\omegaTransitivePart$ is an inductive set.
\end{proposition}
\begin{proof}
    Follows by \auto{inductive_set,omega_transitive_part_empty,omega_transitive_part_successor}.
\end{proof}

\begin{definition}[Natural numbers]\label{omega}
    $\naturals = \{ x \in \omegaUniverse \mid
        \text{for every $A$ such that $A$ is an inductive set
        we have $x\in A$} \}$.
\end{definition}

\begin{proposition}[Zero is natural]\label{omega_empty}
    $\emptyset \in \naturals$.
\end{proposition}
\begin{proof}
    Follows by \auto{omega,omega_universe_empty,inductive_set}.
\end{proof}

\begin{proposition}[Successors are natural]\label{omega_successor}
    For all $x$ we have if $x\in\naturals$, then
    $\suc{x}\in\naturals$.
\end{proposition}
\begin{proof}
    Follows by \auto{omega,infinity,omega_transitive_part,omega_universe_transitive_successor,inductive_set}.
\end{proof}

\begin{proposition}[The natural numbers are inductive]\label{omega_inductive}
    $\naturals$ is an inductive set.
\end{proposition}
\begin{proof}
    Follows by \auto{inductive_set,omega_empty,omega_successor}.
\end{proof}

\begin{proposition}[Inductive-set minimality]\label{omega_minimal}
    For all $A$ we have if $A$ is an inductive set, then
    for all $x$ if $x\in\naturals$, then $x\in A$.
\end{proposition}
\begin{proof}
    Follows by \auto{omega}.
\end{proof}

\begin{abbreviation}[Natural numbers]\label{natural_number}
    $n$ is a natural number iff $n\in\naturals$.
\end{abbreviation}
